\chapter*{Appendix M: Geometric Operators}
\addcontentsline{toc}{chapter}{Appendix M: Geometric Operators}

\section{Metric and Volume Structure}

\subsection{Riemannian Metric}

Let $(\mathcal{M},g)$ be a smooth Riemannian manifold.
The metric is a $(0,2)$-tensor field
\[
g: T\mathcal{M} \times T\mathcal{M} \to \mathbb{R},
\qquad
(v,w) \mapsto g(v,w),
\]
smooth, symmetric, and positive definite.

\subsection{Volume Form}

The Riemannian volume form is
\[
dV_g = \sqrt{\det g_{ij}}\, dx^1\wedge\cdots\wedge dx^n.
\]

\section{Covariant Derivatives}

\subsection{Levi–Civita Connection}

The Levi–Civita connection $\nabla$ is the unique torsion-free connection satisfying metric compatibility:
\[
\nabla_X g = 0.
\]
Christoffel symbols are
\[
\Gamma^k_{ij}
  = \frac12 g^{kl}
    \left( \partial_i g_{jl}
          + \partial_j g_{il}
          - \partial_l g_{ij} \right).
\]

\subsection{Covariant Derivative of Tensors}

For a vector field $X$ and tensor $T$,
\[
(\nabla_X T)^{i_1\cdots i_p}_{j_1\cdots j_q}
= X^k \partial_k T^{i_1\cdots i_p}_{j_1\cdots j_q}
  + \sum_{r=1}^p \Gamma^{i_r}_{kl} X^k
        T^{i_1\cdots l\cdots i_p}_{j_1\cdots j_q}
  - \sum_{s=1}^q \Gamma^l_{kj_s} X^k
        T^{i_1\cdots i_p}_{j_1\cdots l\cdots j_q}.
\]

\section{Divergence, Laplacian, and Gradient}

\subsection{Gradient}

For a smooth scalar field $\phi$, the gradient is the vector field:
\[
\nabla \phi = g^{ij} (\partial_j \phi)\, \partial_i.
\]

\subsection{Divergence}

For a vector field $X$,
\[
\nabla \cdot X = \frac{1}{\sqrt{\det g}}
  \partial_i \bigl(\sqrt{\det g}\, X^i \bigr).
\]

\subsection{Laplace–Beltrami Operator}

For a scalar field $\phi$,
\[
\Delta \phi = \nabla \cdot (\nabla \phi)
  = \frac{1}{\sqrt{\det g}}
    \partial_i\bigl(\sqrt{\det g}\, g^{ij}\partial_j \phi \bigr).
\]

\section{Curvature Operators}

\subsection{Riemann Curvature Tensor}

\[
R(X,Y)Z = \nabla_X\nabla_Y Z - \nabla_Y\nabla_X Z - \nabla_{[X,Y]} Z.
\]

In coordinates:
\[
R^m_{\;\;ijk}
= \partial_i\Gamma^m_{jk}
  - \partial_j\Gamma^m_{ik}
  + \Gamma^m_{il}\Gamma^l_{jk}
  - \Gamma^m_{jl}\Gamma^l_{ik}.
\]

\subsection{Ricci Curvature}

\[
\mathrm{Ric}_{ij} = R^k_{\;\;ikj}.
\]

\subsection{Scalar Curvature}

\[
\mathcal{R} = g^{ij}\mathrm{Ric}_{ij}.
\]

\section{Lie Derivatives}

\subsection{Lie Derivative of a Scalar}

For scalar field $\phi$ along vector $X$:
\[
\mathcal{L}_X \phi = X(\phi).
\]

\subsection{Lie Derivative of a Vector Field}

\[
\mathcal{L}_X Y = [X,Y].
\]

\subsection{Lie Derivative of the Metric}

\[
(\mathcal{L}_X g)_{ij}
  = X^k\partial_k g_{ij}
    + (\partial_i X^k) g_{kj}
    + (\partial_j X^k) g_{ik}.
\]

\section{Hodge Operators}

\subsection{Exterior Derivative}

For a $k$-form $\omega$,
\[
d\omega = \sum_{i_0 < \cdots < i_k}
         \partial_{i_0}\omega_{i_1\cdots i_k}
         \, dx^{i_0}\wedge\cdots\wedge dx^{i_k}.
\]

\subsection{Hodge Star}

The Hodge star $*: \Omega^k(\mathcal{M}) \to \Omega^{n-k}(\mathcal{M})$ satisfies:
\[
\alpha \wedge (\ast \beta)
  = g(\alpha,\beta)\, dV_g.
\]

\subsection{Codifferential}

The adjoint of $d$ is
\[
\delta = (-1)^{nk + n + 1} \ast d \ast.
\]

\subsection{Hodge Laplacian}

\[
\Delta_H = d\delta + \delta d.
\]

\section{Geodesic Operators}

\subsection{Geodesic Equation}

A curve $\gamma(t)$ is geodesic if
\[
\frac{d^2\gamma^k}{dt^2}
  + \Gamma^k_{ij}
    \frac{d\gamma^i}{dt}
    \frac{d\gamma^j}{dt}
  = 0.
\]

\subsection{Geodesic Flow}

The geodesic vector field on $T\mathcal{M}$ is
\[
G = v^i \frac{\partial}{\partial x^i}
  - \Gamma^i_{jk} v^j v^k
    \frac{\partial}{\partial v^i}.
\]

\section{Energy and Action Operators}

\subsection{Dirichlet Energy}

For scalar field $\phi$:
\[
E[\phi]
  = \frac12 \int_{\mathcal{M}}
      g^{ij}(\partial_i\phi)(\partial_j\phi)\, dV_g.
\]

\subsection{Field Action}

A generic field action is
\[
S[\phi] = \int_{\mathcal{M}}
  \mathcal{L}(\phi,\nabla\phi,g)\, dV_g.
\]

\subsection{Euler–Lagrange Operator}

The functional derivative satisfies:
\[
\frac{\delta S}{\delta \phi}
  = \frac{\partial\mathcal{L}}{\partial\phi}
  - \nabla_i \left(\frac{\partial\mathcal{L}}{\partial(\nabla_i\phi)}\right).
\]


